%!TEX root = ../../exa-ma-d7.1.tex
\chapter{Training and capacity building}
\label{chap:training}

\section{Mandate and integration of D7.3}
This chapter integrates Deliverable D7.3, \emph{Training}, v1.0.0, submitted on
15 April 2025, into the Exa-MA mega-deliverable. D7.3 documented
the initial Exa-MA strategy, academic-program survey, 2024 initiatives,
software-training inventory and roadmap for WP7 Task 7.4. The submitted
deliverable is available at
\href{https://drive.google.com/file/d/1d4AKTGNbD74n90d5RESHytOr-M56sFIm}{exa-ma-d7.3-v1.0.0.pdf}.

The v2.90.0 integration draft consolidates rather than reproduces the entire
48-page D7.3. It retains traceability, records what is actually evidenced, and
creates the structure in which WP7 and software owners must report changes since
the version submitted in D7.3.
Plans from D7.3 are never counted as completed activities without dated evidence.

\section{Training strategy and target audiences}
D7.3 defined five audiences: undergraduate and Master's students, doctoral
researchers, HPC engineers and developers, industrial R\&D teams, and academic
or research institutions.  The delivery portfolio combines courses, tutorials,
workshops, schools, webinars, projects, internships, documentation and notebooks.
For v3, each resource should state prerequisite level and connect practical work
to a versioned Exa-MA software, application or benchmark whenever applicable.

\section{Master's and doctoral education connected to Exa-MA}
D7.3 surveyed programmes at the University of Strasbourg,
Grenoble INP/UGA, Nantes Universit\'e, Toulouse INP, Universit\'e Paris-Saclay,
Universit\'e de Lille and UPPA, and documented 2024 student projects and
internships. These records describe what was reported in D7.3; they are not
evidence that the same
courses or projects ran in 2025--2026.

\todo{WP7 Task 7.4 owner: validate the current programme list; for each 2025--2026
course, project, internship or doctoral school provide dates, organizer, WP and
software/application mapping, participant count, materials, outcomes and evidence.}

\section{Schools, workshops, tutorials, hackathons and hands-on sessions}
Only the following completed initiatives were sufficiently described in the
sources inspected for this integration pass.  The normalized records, including
all unknown fields, are in \nolinkurl{data/training-initiatives-v3.csv}.

\begingroup
\footnotesize
\setlength{\tabcolsep}{2pt}
\begin{longtable}{@{}>{\raggedright\arraybackslash}p{0.23\textwidth}>{\raggedright\arraybackslash}p{0.24\textwidth}>{\raggedright\arraybackslash}p{0.16\textwidth}>{\raggedright\arraybackslash}p{0.21\textwidth}>{\raggedright\arraybackslash}p{0.10\textwidth}@{}}
\caption{Verified historical training initiatives inherited from D7.3.}
\label{tab:v3-training-initiatives}\\
\toprule
\textbf{Initiative ID} & \textbf{Title} & \textbf{Date/location} & \textbf{Audience/topic} & \textbf{Status} \\
\midrule
\endfirsthead
\toprule
\textbf{Initiative ID} & \textbf{Title} & \textbf{Date/location} & \textbf{Audience/topic} & \textbf{Status} \\
\midrule
\endhead
\nolinkurl{TR-2024-HIDALGO2-01} & Exa-MA and HiDALGO2 Masterclass HPC & 11 Sep 2024, Strasbourg & Master/PhD; urban-energy exascale overview & completed \\
\nolinkurl{TR-2024-WP7-WEBINAR-01} & Exa-MA WP7 Webinar & 18--19 Nov 2024, Zoom & Researchers/HPC staff; Guix-HPC, Spack and benchmark deployment & completed \\
\bottomrule
\end{longtable}
\endgroup

\todo{Supply the complete 2025--2026 event inventory.  Planned activities must
retain a planned status until dated attendance/material evidence is provided.}

\section{Software-specific training material}
The D7.3 inventory contained the nine software packages below.  The live
\emph{Frameworks} tab confirms that several documentation/training resources are
still registered, but it does not provide dated 2025--2026 initiative or usage
evidence.  No additional software is added in this draft.

\begin{longtable}{@{}>{\raggedright\arraybackslash}p{0.16\textwidth}>{\raggedright\arraybackslash}p{0.34\textwidth}>{\raggedright\arraybackslash}p{0.41\textwidth}@{}}
\caption{D7.3 software-training inventory and v3 validation need.}
\label{tab:v3-training-software}\\
\toprule
\textbf{Software} & \textbf{Baseline resource} & \textbf{v3 status/action} \\
\midrule
\endfirsthead
\toprule
\textbf{Software} & \textbf{Baseline resource} & \textbf{v3 status/action} \\
\midrule
\endhead
CGAL & \url{https://doc.cgal.org/latest/Manual/tutorials.html} & Registered tutorial resource; owner must identify 2025--2026 updates and reuse evidence. \\
Composyx & \url{https://composyx.gitlabpages.inria.fr/composyx/} & Documentation registered; dated training initiative and outcome evidence missing. \\
Feel++ & \url{https://docs.feelpp.org} & Documentation/initial-training resources registered; versioned 2025--2026 modules and KPI data requested. \\
FreeFEM++ & \url{https://doc.freefem.org/introduction/index.html} & Documentation and FreeFEM Days material registered; Exa-MA-specific 2025--2026 reuse must be identified. \\
Hawen & \url{https://ffaucher.gitlab.io/hawen-website/} & Documentation/tutorial resource registered; completion of planned GPU training is not evidenced. \\
HPDDM & \url{https://github.com/hpddm/hpddm/blob/main/doc/cheatsheet.pdf} & Documentation/tutorial resource registered; dated initiative and participant evidence missing. \\
pBB & \url{https://gitlab.inria.fr/jgmys/permutationbb} & Recorded in D7.3 only; current training material and initiative evidence requested. \\
Samurai & \url{https://hpc-math-samurai.readthedocs.io/en/latest/tutorial/} & Tutorial recorded in D7.3; notebook/workshop and benchmark-training roadmap requires validation. \\
Scimba & \url{https://gitlab.inria.fr/scimba/scimba} & D7.3 examples/documentation baseline; API rewrite and multi-GPU training plans require validation. \\
\bottomrule
\end{longtable}

\section{Application- and benchmark-oriented training}
Application training should use the IDs in \Cref{tab:v3-application-registry};
benchmark training should use the evidence stages in
\Cref{tab:v3-benchmark-registry}.  A training exercise must not describe a
proposed benchmark as executable.  The inspected apps-feelpp repository provides
notebooks and command-line recipes, including benchmark execution recipes for
the distance and I/O applications, which supports a \emph{partially achieved}
classification for mini-app and benchmark tutorials.  Participant use, learning
outcomes and reusable releases remain unverified.

\section{Training resources and FAIR/reproducible access}
For each initiative, the editorial registry requires a durable material URL,
software version or commit, license, prerequisites, runnable environment,
application/benchmark ID and evidence of the delivered session.  Where data are
needed, the benchmark registry's configuration, machine and raw-data fields
apply.  A central portal has not been verified; current resources are distributed
across software documentation sites and repositories.

\section{Outcomes and KPIs}
D7.3 proposed attendance, repository analytics and feedback surveys as KPI
sources.  The two historical events above do not include participant counts,
language, material URLs or feedback/reuse indicators in the submitted report.
No aggregate 2025--2026 attendance or feedback dataset was supplied during this
integration pass, so no new KPI value is reported.

\todo{Task 7.4 owner: provide a deduplicated initiative export with attendance,
audience, language, materials, feedback/reuse indicator and consent-compliant
evidence; reconcile it with the project KPI workbook.}

\section{Gaps and roadmap for 2027--2029}
\begin{longtable}{@{}>{\raggedright\arraybackslash}p{0.38\textwidth}>{\raggedright\arraybackslash}p{0.18\textwidth}>{\raggedright\arraybackslash}p{0.35\textwidth}@{}}
\caption{Conservative review of objectives inherited from the D7.3 roadmap.}
\label{tab:v3-training-roadmap}\\
\toprule
\textbf{D7.3 objective} & \textbf{Classification} & \textbf{Evidence gap / next action} \\
\midrule
\endfirsthead
\toprule
\textbf{D7.3 objective} & \textbf{Classification} & \textbf{Evidence gap / next action} \\
\midrule
\endhead
Upgrade tutorials with GPU, Kokkos or multi-node examples & not started & No dated training release or delivered initiative was supplied; owners must link versioned modules. \\
Create mini-app training modules & partially achieved & Interactive notebooks exist in the apps-feelpp repository; identify releases, audiences and actual reuse. \\
Create benchmark tutorials & partially achieved & Distance and I/O execution recipes exist; package them with pinned machines/data and learning outcomes. \\
Document software-development and reproducibility guidance & partially achieved & Repository and contribution guidance exists; consolidate it into a maintained training module. \\
Provide a centralized training portal & ongoing & Distributed documentation exists; nominate portal owner, catalogue and preservation policy. \\
Collect attendance and feedback systematically & not started & No 2025--2026 consolidated dataset was supplied; adopt the normalized initiative schema. \\
Collaborate with NumPEx PCs, GENCI and computing centres & not started & No dated post-D7.3 collaboration evidence was supplied; record partners, event and evidence before claiming completion. \\
\bottomrule
\end{longtable}

For 2027--2029, priority should be given to closing evidence gaps: publish
versioned learning packages around executable applications/benchmarks, capture
attendance and feedback consistently, and preserve materials and environments in
durable repositories.  Scientific or scheduling commitments remain for WP7 and
the software owners to approve.
